% ==========================================================================
%  introduction.tex — Physically-Guided Collision-Free Indoor Scene Generation
% ==========================================================================
\section{Introduction}
\label{sec:intro}

Synthesizing realistic 3D indoor scenes is a fundamental problem in computer vision and graphics, with direct applications in virtual reality~\cite{placeholder_vr}, robotics simulation~\cite{placeholder_robotics}, interior design~\cite{placeholder_interior}, and training data generation for embodied AI~\cite{placeholder_embodied}.
A practical scene generation system must satisfy two requirements simultaneously:
\textbf{(i)}~semantic plausibility---objects should be arranged in configurations that respect high-level spatial relationships (e.g., a nightstand \emph{beside} a bed, chairs \emph{surrounding} a table); and
\textbf{(ii)}~physical validity---objects must not interpenetrate, must remain within room boundaries, and the resulting layout should be navigable by a human or robotic agent.

Recent advances in scene graph-conditioned generation have made significant progress toward the first requirement.
By representing a scene as a graph $\mathcal{G} = (\mathcal{V}, \mathcal{E})$ whose nodes encode object categories and edges encode spatial predicates (e.g., \texttt{left-of}, \texttt{close-by}), methods such as Graph-to-3D~\cite{placeholder_graph2threed}, CommonScenes~\cite{placeholder_commonscenes}, and EchoScene~\cite{echoscene} can generate diverse, controllable layouts that faithfully follow the input specification.
EchoScene, in particular, introduces a dual-branch diffusion framework with an \emph{Information Exchange Unit} that propagates each object's denoising state to its graph neighbors at every reverse step, achieving state-of-the-art fidelity and manipulation robustness.

However, \textbf{physical validity remains an open challenge}.
The training objectives of existing scene graph diffusion models---typically reconstruction losses over bounding box parameters---do not explicitly penalize inter-object collisions.
The Information Exchange Unit in EchoScene provides implicit relational awareness that helps objects distribute spatially, but it offers no hard physical guarantees.
As illustrated in Fig.~\ref{fig:teaser}, generated scenes frequently exhibit interpenetrating furniture, objects protruding outside room walls, and cluttered central areas that would be impassable in practice.
These artifacts severely limit the applicability of generated scenes for downstream tasks such as robotic navigation planning~\cite{placeholder_nav_planning} or physics-based simulation~\cite{placeholder_physics_sim}, where even minor overlaps cause solver failures.

A natural question arises: \emph{how can we enforce physical constraints without sacrificing the semantic coherence learned by the diffusion model?}
One approach is to add collision penalties during training, but this requires costly retraining and risks destabilizing the learned distribution.
Another is to apply heuristic post-processing to push apart overlapping objects, but large post-hoc displacements can destroy the carefully generated spatial relationships.
We observe that neither approach alone is sufficient: training-time losses lack geometric precision at inference, while post-processing lacks awareness of semantic context.

In this paper, we propose a \textbf{coarse-to-fine collision resolution framework} that combines the best of both paradigms.
Our approach operates in two complementary stages:

\begin{enumerate}[nosep]
    \item \textbf{Physics-Guided Diffusion (Coarse Stage).}
    During the reverse denoising process, we inject three differentiable guidance losses---a collision loss based on 3D oriented IoU and penetration volume, a room-layout boundary loss, and a reachability loss---that steer the predicted clean layout $\hat{\mathbf{d}}_0^{(t)}$ toward physically valid configurations at each timestep.
    Crucially, this guidance is \emph{training-free}: we compute loss gradients with respect to the predicted layout and use them to shift the posterior mean, following the classifier-guidance paradigm~\cite{placeholder_classifier_guidance}.
    The model's pretrained weights remain frozen, preserving its learned distribution while redirecting samples away from collisions.

    \item \textbf{OBB-SAT Post-Processing (Fine Stage).}
    After the full reverse process completes, we apply a deterministic collision resolution module that uses the Separating Axis Theorem (SAT) on Oriented Bounding Boxes (OBBs) in the floor plane.
    This stage detects and resolves any residual interpenetrations with geometrically exact, rotation-aware checks, making only small, local corrections since the coarse stage has already handled the bulk of the overlaps.
\end{enumerate}

This decomposition is effective because the two stages address complementary failure modes.
The coarse guidance operates in the model's learned latent space during generation, maintaining semantic coherence while broadly reducing collisions.
The fine post-processing operates on the final metric-space output with exact oriented geometry, eliminating residual overlaps that the approximate gradient-based guidance cannot fully resolve.
Together, they achieve collision-free layouts without the large displacements that degrade semantic quality.

\paragraph{Contributions.}
Our main contributions are as follows:

\begin{itemize}[nosep]
    \item We propose a \textbf{training-free physical guidance} mechanism for scene graph layout diffusion that integrates three complementary losses---collision avoidance, room-boundary enforcement, and reachability---into the reverse denoising process via variance-preconditioned gradient steering.

    \item We introduce an \textbf{OBB-SAT post-processing} module that performs rotation-aware collision detection and center-biased resolution in the floor plane, providing geometrically exact fine-grained correction.

    \item We demonstrate that the \textbf{coarse-to-fine combination} of both stages significantly reduces inter-object collisions compared to the EchoScene baseline and post-processing-only approaches, while preserving or improving layout quality metrics on the SG-FRONT benchmark~\cite{placeholder_sgfront}.
\end{itemize}
